\section{Conclusions \& Reflections}
\label{sec:conclusion}

This project successfully achieved all the requirements of the coursework, and constructed a complete and stable Raspberry Pi software and hardware interaction system. In terms of basic functions, the program directly controls the physical keys, LED indicators and display through the underlying code, and properly handles the timeout problem when the user input. In the core cracking logic, the system not only uses efficient basic assembly instructions to complete the comparison calculation of password differences, but also implements an intelligent derivation algorithm for arbitrary length passwords. All expected tasks have passed rigorous system testing.

The highlight of our project lies in its system optimization and integration of hardware and software. In terms of hardware, fast response is achieved by directly manipulating memory addresses. As for algorithm, the system uses logical investigation and advance filtering instead of blind brute force testing, which fundamentally avoids the computing power collapse that may be caused when long passwords are cracked. At the same time, the program insists on applying memory to the system on demand, and strictly controls the resource occupation during the running process, which perfectly adapts to the objective reality of the limited hardware capabilities of embedded devices.

Through this system programming practice, we have a deeper understanding of the software and hardware co-development. It significantly helped to improve our ability to write low-level instructions and learn how to manage memory more safely and precisely in C. More importantly, this assignment helped us establish a global system thinking, that is, we learned how to transform the complex software logic into efficient and stable bottom execution actions step by step with limited physical resources, which laid a solid foundation for the future engineering development.
